\begin{frame}{Az algoritmus alapötlete 1.}
  \small
  \begin{itemize}
%    \item Given T and the remaining non-tree edges E\\ET sorted from lowest to highest weight.
    \item Adott T és a nem-fa élek $E \setminus E_T$ súly szerint növekvő sorban.
%    \item Observe that each of the m - n + 1 edges in E\\ET induces a fundamental cycle with the edges in T.
    \item Figyeljük meg, hogy mindegyik m - n + 1 darab él $E \setminus E_T$-ben egy fundamentális kört indukál T éleivel.
%    \item Then for any MST edge there is a subset of cycles containing that edge, and the cycle induced by the lightest non-MST edge is the replacement for it
    \item Then for any MST edge there is a subset of cycles containing that edge, and the cycle induced by the lightest non-MST edge is the replacement for it
  \end{itemize}
  \centering
  \begin{tikzpicture}[
    scale=1,
    vtx/.style={circle, draw=black, fill=white, inner sep=1pt, minimum size=16pt},
    mstedge/.style={line width=2.5pt, draw=black},
    otheredge/.style={line width=0.6pt, draw=black},
    font=\small
  ]

    % --- Node coordinates ---
    \node[vtx] (a) at (-4.2, 1) {a};
    \node[vtx] (b) at (-2.2, 2) {b};
    \node[vtx] (d) at (-3, -1.2) {d};
    \node[vtx] (e) at (0, 1) {e};
    \node[vtx] (f) at (-0.5, -1.2) {f};
    \node[vtx] (g) at (2, 2) {g};
    \node[vtx] (h) at (2.5, 0) {h};

    % -- root node
    \node[rootvertex] (c) at (-2.0, 0.2) {c};

    % --- MST edges (thick) ---
    \draw[mstedge] (d) -- (f)
      node[midway, below=2pt, fill=white, inner sep=1pt, font=\scriptsize]{1};
    \draw[mstedge] (g) -- (h)
      node[midway, right=2pt, fill=white, inner sep=1pt, font=\scriptsize]{3};
    \draw[mstedge] (e) -- (f)
      node[midway, right=2pt, fill=white, inner sep=1pt, font=\scriptsize]{4};
    \draw[mstedge] (b) -- (e)
      node[midway, below=2pt, fill=white, inner sep=1pt, font=\scriptsize]{5};
    \draw[mstedge] (c) -- (d)
      node[midway, left=3pt, fill=white, inner sep=1pt, font=\scriptsize]{7};
    \draw[mstedge] (a) -- (d)
      node[midway, left=3pt, fill=white, inner sep=1pt, font=\scriptsize]{9};
    \draw[mstedge] (e) -- (h)
      node[midway, above=2pt, fill=white, inner sep=1pt, font=\scriptsize]{2};

    % --- Non-tree edges (thin) ---
    \draw[otheredge] (b) -- (g)
      node[midway, above=2pt, fill=white, inner sep=1pt, font=\scriptsize]{8};
    \draw[otheredge] (a) -- (b)
      node[midway, above=3pt, fill=white, inner sep=1pt, font=\scriptsize]{10};
    \draw[otheredge] (c) -- (e)
      node[midway, above=2pt, fill=white, inner sep=1pt, font=\scriptsize]{11};
    \draw[otheredge] (b) -- (c)
      node[midway, left=2pt, fill=white, inner sep=1pt, font=\scriptsize]{12};
    \draw[otheredge] (f) -- (h)
      node[midway, above=2pt, fill=white, inner sep=1pt, font=\scriptsize]{13};
    \draw[otheredge] (e) -- (g)
      node[midway, above=2pt, fill=white, inner sep=1pt, font=\scriptsize]{6};

  \end{tikzpicture}

  \medskip
  \small Itt: az $(e,g)$ él az $(e,h)$ és a $(g,h)$ élek minimális költségű csereéle.
\end{frame}

\begin{frame}{Az algoritmus alapötlete 2.}
  \small
  \begin{enumerate}
    \item Assign the parent and the first and last visited numbers to each vertex according to depth-first search over T
    \item Traverse the fundamental cycle induced by each non-tree edge in order of ascending weight
    \item Compress paths to skip edges already assigned replacements
  \end{enumerate}
  We use the disjoint sets for fast path compression based on the Gabow-Tarjan static union tree method.


\centering
\begin{tikzpicture}[
  scale=1,
  vtx/.style={circle, draw=black, fill=white, inner sep=1pt, minimum size=16pt},
  mstedge/.style={line width=2.5pt, draw=black},
  otheredge/.style={line width=0.6pt, draw=black},
  font=\small
]

% --- Node coordinates ---
\node[vtx] (a) at (-4.2, 1) {a};
\node[vtx] (b) at (-2.2, 2) {b};
\node[vtx] (d) at (-3, -1.2) {d};
\node[vtx] (e) at (0, 1) {e};
\node[vtx] (f) at (-0.5, -1.2) {f};
\node[vtx] (g) at (2, 2) {g};
\node[vtx] (h) at (2.5, 0) {h};

% -- root node
\node[rootvertex] (c) at (-2.0, 0.2) {c};

% --- MST edges (thick) ---
\draw[mstedge] (d) -- (f)
  node[midway, below=2pt, fill=white, inner sep=1pt, font=\scriptsize]{1};
\draw[mstedge] (g) -- (h)
  node[midway, right=2pt, fill=white, inner sep=1pt, font=\scriptsize]{3};
\draw[mstedge] (e) -- (f)
  node[midway, right=2pt, fill=white, inner sep=1pt, font=\scriptsize]{4};
\draw[mstedge] (b) -- (e)
  node[midway, below=2pt, fill=white, inner sep=1pt, font=\scriptsize]{5};
\draw[mstedge] (c) -- (d)
  node[midway, left=3pt, fill=white, inner sep=1pt, font=\scriptsize]{7};
\draw[mstedge] (a) -- (d)
  node[midway, left=3pt, fill=white, inner sep=1pt, font=\scriptsize]{9};
\draw[mstedge] (e) -- (h)
  node[midway, above=2pt, fill=white, inner sep=1pt, font=\scriptsize]{2};

% --- Non-tree edges (thin) ---
\draw[otheredge] (b) -- (g)
  node[midway, above=2pt, fill=white, inner sep=1pt, font=\scriptsize]{8};
\draw[otheredge] (a) -- (b)
  node[midway, above=3pt, fill=white, inner sep=1pt, font=\scriptsize]{10};
\draw[otheredge] (c) -- (e)
  node[midway, above=2pt, fill=white, inner sep=1pt, font=\scriptsize]{11};
\draw[otheredge] (b) -- (c)
  node[midway, left=2pt, fill=white, inner sep=1pt, font=\scriptsize]{12};
\draw[otheredge] (f) -- (h)
  node[midway, above=2pt, fill=white, inner sep=1pt, font=\scriptsize]{13};
\draw[otheredge] (e) -- (g)
  node[midway, above=2pt, fill=white, inner sep=1pt, font=\scriptsize]{6};

% --- DFS (IN, OUT) pairs ---
\node[font=\scriptsize] at ($(a)+(-0.2,0.5)$) {(13, 14)};
\node[font=\scriptsize] at ($(b)+(0,0.6)$) {(9, 10)};
\node[font=\scriptsize] at ($(c)+(-0.6,0.4)$) {(1, 16)};
\node[font=\scriptsize] at ($(d)+(-0.1,-0.6)$) {(2, 15)};
\node[font=\scriptsize] at ($(e)+(0,0.5)$) {(4, 11)};
\node[font=\scriptsize] at ($(f)+(0.0,-0.6)$) {(3, 12)};
\node[font=\scriptsize] at ($(g)+(0,0.6)$) {(6, 7)};
\node[font=\scriptsize] at ($(h)+(0.4,0.5)$) {(5, 8)};

\end{tikzpicture}
\end{frame}